\origpage{23}{GRAPH MINORS. XIV}
{\centering
{\Large Graph Minors\par}
{\large XIV. Extending an Embedding\par}\medskip

Neil Robertson$^*$\par
\textit{Department of Mathematics, Ohio State University, 231 West 18th Avenue, Columbus, Ohio 43210}\par
AND\par
P. D. Seymour\par
\textit{Bellcore, 445 South Street, Morristown, New Jersey 07960}\par
Received March 15, 1991\par
}
\begin{quote}
Let $G$ be a graph, and let $H$ be a subgraph of $G$ drawn in a surface $\Sigma$. When can this drawing be extended to an embedding of the whole of $G$ in $\Sigma$, up to 3-separations? We show that if such an extension is impossible, and if $H$ is a subdivision of a simple 3-connected graph and is highly ``representative'', then one of two obstructions is present. This is a lemma for use in a future paper.
\end{quote}

\sectionhead{1. Introduction}

The next goal of this series of papers is to show that, for any fixed graph $H$, all graphs with no minor isomorphic to $H$ are tree-structures of pieces which in some sense are ``almost'' of low genus. We are therefore concerned with theorems saying that graphs not possessing certain obstructions can be embedded in simple surfaces. One natural question is, given a graph $G$ and a surface $\Sigma$, when can $G$ be embedded in $\Sigma$? But answering this in a way which is useful for proving theorems (rather than an algorithmic answer) is probably close to finding explicitly all the excluded minors for embedding in $\Sigma$, and is therefore much too difficult. Fortunately, we do not really need to answer this question. In our application, we already have a subgraph $H$ of $G$ drawn in $\Sigma$ which is highly ``representative'' (that is, every non-null-homotopic closed curve in $\Sigma$ meets the drawing of $H$ a large number of times), and it suffices for us to find a condition under which this%
\origfootnote{Research partially supported by NSF Grant DMS8504054.}
